\scanpage{117}



\deleted{4}\handwritten{5}.6. Variantes avec bi-proalgèbres .

\deleted{4}\handwritten{5}.7. Variantes graduées.

\uline{Par}.\deleted{5}\handwritten{6} . Bigèbres d'endomorphismes des $\otimes$-foncteurs. Schémas en groupes d'automorphismes de $\otimes$-foncteurs.

\figure{\begin{tikzpicture}[x=1em,y=1em,baseline]
\node at (0,0) {.};\node at (.7,0) {.};\node at (1.4,0) {.};\node at (2.1,0) {.};
\draw[line width=.15pt] (-.3,.35)--(-.3,.65)--(-.1,.63);
\draw[line width=.15pt] (.05,.63)--(.08,.35);
\draw[line width=.15pt] (.27,.35) .. controls (.45,.43) and (.46,.2) .. (.42,.15);
\end{tikzpicture}}

\uline{Par} \deleted{6}\handwritten{7}. Structure des catégories de Picard\handwritten{(?)}(Vi\unsure{\deleted{n}e\deleted{e}n|\deleted{e}n\deleted{n}e|nn}drait mieux avant par.4).

\margin{\handwritten{N'est pas \unsure{nécessairement|necessairement} le lieu ici}}

\handwritten{\deleted{\uline{Par}.8 \unsure{É|E!}} Plus gén, \deleted{\unsure{catégories}}\inserted{groupoïdes} avec $\otimes$ associatif et unitaire (pas néc. comm. ...).}

\margin{\handwritten{resp. com.}}
